\mypart{Пределы}

\section{Топология и метрика}

Для дальнейшего нам потребуется понятие расстояния между точками. Пусть~$A$~"--- некоторое множество.

\subsection{Открытые множества и расстояния}

\begin{definition}
	Будем называть функцию~$\rho \colon A \times A \to [0, +\infty)$
	\tdef{расстоянием}, или \tdef{метрикой}, если для любых~$x$, $y$, $z \in A$
	она удовлетворяет следующим свойствам:%
	\par
	\addvspace{\medskipamount}%
	\begin{tabular}{rl@{\qquad}l}
		1) & $\rho(x, y) \geq 0$ & (неотрицательность);
		\\[\smallskipamount]
		2) & $\bigl( \rho(x, y) = 0 \bigr) \iff \bigl( x = y \bigr)$ &
			 (невырожденность);
		\\[\smallskipamount]
		3) & $\rho(x, y) = \rho(y, x)$ & (симметричность);
		\\[\smallskipamount]
		4) & $\rho(x, z) \leq \rho(x, y) + \rho(y, z)$ & (неравенство треугольника).
	\end{tabular}
\end{definition}

\begin{example}
	Пример расстояния между точками в~$\RR$~"--- модуль: если $a$ и~$b$~"--- две
	точки из~$\RR$, то можно определить расстояние~$\rho_1$ формулой
	\[
		\rho_1(a, b) := \abs{a - b}.
	\]
	Почему это метрика? Проверим, выполняются ли условия 1)--4) из определения
	метрики:
	\begin{enumerate}[1)]
		\item $\rho_1(a, b) = \abs{a - b} \geq 0$ по определению модуля
			вещественного числа;
		\item если $\rho_1(a, b) = \abs{a - b} = 0$, то $a = b$; обратно, если $a =
			b$, то $\rho_1(a, b) = \rho_1(a, a) = \abs{a - a} = 0$;
		\item $\rho_1(a, b) = \abs{a - b} = \abs{-(b - a)} = \abs{b - a} =
			\rho_1(b, a)$;
		\item $\rho_1(a, c) = \abs{a - c}$ и $\rho_1(a, b) + \rho_1(b, c) = \abs{a
			- b} + \abs{b - c}$, и, как известно%
			\footnote{%
				Тут нужно вспомнить, как раскрываются модули, и вывести отсюда факт:
				\[
					\abs{x + y} \leq \abs{x} + \abs{y}
					.
				\]
			}%
			,
			\[
				\abs{a - c} =
				\abs{a - b + b - c} \leq
				\abs{a - b} + \abs{b - c}
				.
			\]
	\end{enumerate}
\end{example}

\subsection{Окрестности и шары}

\begin{definition}
	Подмножество~$A \subset B$ называется \tdef{открытым в~$B$}, если вместе с
	каждой своей точкой оно содержит шар (положительного радиуса):
	\[
		(\text{$A$ открыто в~$B$}) :=
		\left( 
			\forall a \in A\ 
			\exists \epsilon > 0
			\suchthat
			B_{\epsilon}(a) \subset A
		\right)
		.
	\]
\end{definition}

Пусть у нас есть множество~$A$.

\begin{definition}
	\tdef{Окрестность $U(a)$} элемента~$a$ множества~$A$~"--- это открытое
	подмножество множества~$A$, содержащее точку~$a$.
\end{definition}

Если мы умеем измерять расстояние между точками в~$A$, то приобретают смысл следующие определения.
\begin{definition}
	\tdef{Шар~$B_r(a)$ радиуса~$r$ с центром в точке~$a$} (иногда ещё говорят
	\tdef{открытый шар})~"--- это множество таких
	точек из~$A$, что расстояние от~$a$ до них строго меньше~$r$:
	\[
		B_r(a) :=
		\set{x \in A}{\rho(x, a) < r}
		.
	\]
\end{definition}

\begin{definition}
	\tdef{Сфера~$S_r(a)$ радиуса~$r$ с центром в точке~$a$}~"--- это множество
	таких точек из~$A$, что расстояние от~$a$ до каждой из них равно~$r$:
	\[
		S_r(a) :=
		\set{x \in A}{\rho(x, a) = r}
		.
	\]
\end{definition}

\begin{definition}
	\tdef{Проколотый шар~$\Bo_r(a)$ радиуса~$r$ с центром в точке~$a$}~"--- это
	шар~$B_r(a)$, из которого выкинули центр:
	\[
		\Bo_r(a) := B_r(a) \setminus \set*{a} =
		\set{x \in A}{0 < \rho(x, a) < r}
		.
	\]
\end{definition}

В дальнейшем для простоты будем под окрестностью точки понимать шар с центром в
этой точке. Если нужно будет указать радиус~$\epsilon$ шара, то будем говорить
<<$\epsilon$"=окрестность>>, а если шар проколотый~"--- то <<проколотая
$\epsilon$"=окрестность>>.

\section{Предел последовательности}

\begin{definition}
	\tdef{Последовательность} элементов множества~$A$~"--- это отображение~$\phi$
	из~$\NN$ в~$A$. При этом пишут $\phi(m) = a_m$.

	Часто вместо отображения $\phi \colon \NN \to A$ последовательностью называют
	образ этого отображения, то есть множество
	\[
		\phi(\NN) = \set{\phi(n)}{n \in \NN} = \set{a_n}{n \in \NN}
		.
	\]
	Часто пишут просто <<последовательность~$\set*{a_n}$>>, не указывая явно,
	какое множество пробегает индекс~$n$, если это понятно из контекста.
\end{definition}
Иными словами, каждому натуральному числу~$n \in \NN$ сопоставляется какой-то
элемент~$a_n$ множества~$A$. При этом может быть так, что это отображение не
инъективно, то есть разным числам~$m$ и~$n \in \NN$ сопоставляется один и тот
же элемент: $a_m = a_n$ при $m \neq n$.


\begin{definition}
	\tdef{Предел последовательности $\set*{a_n}$}~"--- это точка, в любом открытом
	шаре вокруг которой находятся все элементы этой последовательности, начиная с
	некоторого.
\end{definition}

Формализуем это определение.

Пусть $\phi\colon \NN \to A$~"--- последовательность и~$a \in A$~"--- её предельная точка. Пусть $\epsilon > 0$~"--- произвольное положительное число (по смыслу~"--- радиус шара), а $B_\epsilon(a)$~"--- открытый шар радиуса~$\epsilon$ вокруг точки~$a$:
\begin{align}
	B_\epsilon(a) :=
	{} &
	\set{x \in A}{\rho(x, a) < \epsilon}
	\label{eq:Limits:B_eps_a}
	= {} \\ {} = {} &
	\set{x \in A}{\abs{x - a} < \epsilon}
	.
	\notag
\end{align}
Пусть, далее, $N \in \NN$~"--- натуральный номер элемента, начиная с которого, все элементы последовательности находятся внутри шара~$B_\epsilon(a)$ (и этот номер зависит от радиуса~$\epsilon$). Если это так, то для всех натуральных индексов~$n \in \NN$, как только $n > N$, выполняется услловие $a_n \in B_\epsilon(a)$, то есть (подставим в~условие~\eqref{eq:Limits:B_eps_a}, задающее множество~$B_\epsilon(a)$, вместо~$x$ элемент $a_n$)
\[
	\rho(a_n, a) = \abs{a_n - a} < \epsilon.
\]

Итак, осталось записать всё вместе. Во-первых, мы зафиксировали произвольное положительное число:
\[
	\forall \epsilon > 0
	\hphantom{%
		\ 
		\exists N \in \NN\ 
		\forall n > N
		\holds
		a_n \in B_\epsilon(a)
		.%
	}
	%;
\]
далее, мы сказали, что должен найтись натуральный номер~$N \in \NN$, зависящий от~$\epsilon$:
\[
	\forall \epsilon > 0\ 
	\exists N \in \NN
	\hphantom{%
		\ 
		\forall n > N
		\holds
		a_n \in B_\epsilon(a)
		.%
	}
	%;
\]
дальше говорим, что для всех номеров, больших~$N$, выполняется какое-то выражение (нас пока не интересует, какое именно)~"--- так и запишем:
\[
	\forall \epsilon > 0\ 
	\exists N \in \NN\ 
	\forall n > N
	\hphantom{%
		\holds
		a_n \in B_\epsilon(a)
		.%
	}
	%;
\]
а дальше пишем, что именно происходит при этих условиях:
\[
	\forall \epsilon > 0\ 
	\exists N \in \NN\ 
	\forall n > N
	\holds
	a_n \in B_\epsilon(a)
	.
\]
Или, что то же самое,
\[
	\forall \epsilon > 0\ 
	\exists N \in \NN\ 
	\forall n > N
	\holds
	\abs{a_n - a} < \epsilon
	.
\]

% TODO: написать про отрицание

\section{Предел функции}

Мы уже познакомились с понятием предела последовательности. Определим теперь предел функции. Пусть $f \colon X \to Y$~"--- функция. Эта запись означает, что область определения $f$~"--- всё множество~$X$ целиком и что функция $f$ берёт элемент $x$ множества $X$ и переводит его в какой-то элемент $f(x)$ множества $Y$. Можно считать, что и $X$, и $Y$~"--- это числа, но это не обязательно: $X$ и $Y$ могут иметь произвольную природу.

Для того чтобы определить предел функции, нам понадобятся вспомогательные (важные и сами по себе) понятия.


Определим теперь предел функции.

\begin{definition}
	Элемент $a \in Y$ называется \tdef{пределом функции~$f \colon X \to Y$ при $x
	\to x_0$}, если для любой $\epsilon$-окрестности~$U_\epsilon(a)$
	точки~$a \in Y$ существует проколотая
	$\delta$-окрестность~$\Vo_\delta(x_0)$ точки~$x_0 \in X$, такая, что $f
	\bigl( \Vo_\delta(x_0) \bigr) \subset U_\epsilon(a)$, то есть если любая
	точка из~$\Vo_\delta(x_0)$ (находящаяся на расстоянии не больше чем~$\delta$
	от точки~$x_0$ и не совпадающая с~$x_0$) переходит в какую-то точку
	окрестности~$U_\epsilon(\ell)$, то есть образ $\delta$-окрестности
	$\Vo_\delta(x_0)$ точки~$x_0$ не выходит за пределы $\epsilon$-окрестности
	$U_\epsilon(\ell)$ точки~$\ell$:
	\[
		\left( 
			\lim_{x \to x_0} f(x) = \ell
		\right)
		:=
		\left( 
			\forall U_\epsilon(\ell) \subset Y\ 
			\exists \Vo_\delta(x_0) \subset X\ 
			\suchthat
			f \bigl( \Vo_\delta(x_0) \bigr) \subset
			U_\epsilon(\ell)
		\right)
		.
	\]
\end{definition}

Говоря \emph{очень} неформально, точка~$a \in Y$ называется \tdef{пределом
функции~$f\colon X \to Y$} при $x \to x_0$, если при <<малом шевелении>>
аргумента возле точки~$x_0$ значение $f(x)$ будет <<шевелиться предсказуемо>>.

Формулизуем. Пусть $\epsilon > 0$~"--- радиус окрестности точки~$a \in Y$ (шарика~$B_{\epsilon}(a) \subset Y$), которая формализует понятие <<шевелиться предсказуемо>>

\begin{definition}
\end{definition}

% \ex{1}

